% ---------------------------------------------------------------------
%  Avant-propos
% ---------------------------------------------------------------------
\chapter*{Avant-propos}
\addcontentsline{toc}{chapter}{Avant-propos}
\markboth{Avant-propos}{Avant-propos}

En une quinzaine d'années, la façon dont les organisations stockent et exploitent leurs données a
profondément changé. L'entrepôt de données, conçu pour des tableaux de bord sur des données
structurées, a été complété par le lac de données, capable d'accueillir à faible coût des textes,
des images ou des journaux d'événements, mais sans les garanties de fiabilité d'une base de
données. Le \emph{Lakehouse} réunit ces deux mondes : des fichiers ouverts, stockés à faible coût,
sur lesquels une couche de métadonnées apporte les transactions, l'historique des versions et le
contrôle du schéma. Il est devenu en quelques années l'architecture de référence des plateformes de
données.

Cet ouvrage est né d'un module d'enseignement consacré à cette architecture. Il en reprend la
démarche : comprendre en construisant. Plutôt que de décrire les outils un à un, il fait réaliser au
lecteur une plateforme complète, sur une seule et même entreprise fictive mais réaliste, un site de
commerce en ligne brésilien dont les données ont été rendues publiques. Chapitre après chapitre, la
plateforme s'enrichit : des tables issues d'une base de vente, puis des avis clients et des messages
JSON, des photos de produits, et enfin des événements reçus en continu. À chaque étape, le lecteur
rencontre les questions qu'une équipe de données affronte réellement : la qualité des données
reçues, leur modélisation, l'automatisation des traitements, leur supervision, et le choix entre un
traitement par lots et un traitement en continu.

Les travaux pratiques s'appuient exclusivement sur des logiciels libres, très utilisés en
entreprise : Apache Spark, Apache Iceberg, Delta Lake, dbt, Apache Airflow, Trino, Apache Superset et
Apache Kafka. L'ensemble tourne sur un simple ordinateur portable, dans des conteneurs Docker, et
s'installe en une commande. Le lecteur n'a besoin d'aucun compte chez un fournisseur de cloud, et
les notions acquises se transposent directement aux plateformes commerciales, auxquelles l'ouvrage
fait référence chaque fois que c'est utile.

L'ouvrage s'adresse aux étudiants de master et d'école d'ingénieurs en informatique, en science des
données ou en informatique décisionnelle, ainsi qu'aux ingénieurs qui souhaitent se former à ces
architectures. Il suppose une pratique de base de Python et de SQL, et des notions sur les bases de
données relationnelles.

Nous espérons qu'il permettra au lecteur d'acquérir non seulement la maîtrise de ces outils, mais
aussi le raisonnement qui permet de les choisir, de les combiner et d'en mesurer les limites.

\bigskip
\hfill\textit{Les auteurs}
